\section{Experimental design}
\label{sec:experimental-design}

\subsection{Stable-regime synthetic benchmark}

Real metadata anchors have no physical-bias ground truth. Known-effect
evaluation therefore uses only stable target and donor regimes separated from
nearby reported transitions. The frozen \texttt{stable\_full\_v2} manifest
contains \SyntheticCases{} physical target sites. Whole connected components of
the complete target-plus-donor physical input graph were assigned to the
\CalibrationCases{}-site calibration subset or the \EvaluationCases{}-site
held-out evaluation subset, so no physical input site crosses the split.

\begin{figure}[tbp]
\centering
\includegraphics[width=0.91\linewidth]{fig_split_integrity.pdf}
\caption{Leakage-resistant stable-case split. Whole connected components of
the complete target-plus-donor physical-site graph were assigned to the
calibration or held-out evaluation partition before metrics were computed.
The frozen split audit records zero shared physical input sites. This controls
physical input reuse across partitions; it does not make the synthetic
perturbations a model of every real monitoring process.}
\label{fig:split-integrity}
\end{figure}

The benchmark injects target-only additive, proportional, gradual-drift,
temporary-step, and variance changes, together with matched target-and-donor
regional variants. Each perturbation variant contains \SamplesPerVariant{}
held-out evaluation samples. Each family therefore combines a local condition
and a matched regional condition; the family rows in the results are not
local-only effect estimates. Five frozen magnitude multipliers
(0.5, 1.0, 1.5, 2.0, and 3.0) scale an additive step, proportional step,
gradual drift, temporary step, or variance increase. The variance-only
condition has no defined level-effect truth, so local-effect MAE is correctly
reported as N/A for that condition.

The effect shape, magnitudes, seeds, candidate rule, model settings, and
metrics were frozen before evaluation. For every method, a ranking threshold
is selected only by macro-F1 on the calibration partition and then applied
unchanged to the evaluation partition. The viewed evaluation partition is not
used for model, threshold, donor-weight, quality-gate, or taxonomy adjustment.
This is a protocol safeguard, not a claim that one finite split eliminates all
future generalization uncertainty.

\subsection{Benchmark metrics and uncertainty}

For target-only non-variance perturbations, local-effect MAE is the mean
absolute difference between the estimated and injected 60-day median log
effect. AUPRC measures ranking of local against matched regional variants;
macro-F1, precision, recall, and regional false-positive rate use the
calibration-only threshold. CUSUM, PELT, and rolling-MAD have no
protocol-defined effect-magnitude estimate, so their MAE cells are N/A while
their ranking metrics remain observable.

The MetaShift-versus-standard-MAE comparison uses a paired cluster bootstrap
over stable case IDs with 1,000 repetitions and fixed seed 20260830. Pairing
keeps a method comparison tied to the same pseudo-anchor, perturbation,
magnitude, and injected seed. It reports uncertainty for the frozen benchmark
comparison, rather than a universal probability that one method dominates
another.

\subsection{Real-anchor audit and sensitivity}

The real-anchor workflow uses the same distinct-physical-donor construction,
pre-event fitting boundaries, and effect window across the three common
cross-site methods. The five single-series baselines use the same 60-day
before/after window but do not receive a geographic counterfactual. The audit
records complete comparisons and explicit non-complete dispositions for all
\TotalAnchors{} anchors. Cross-validated MetaShift is evaluated only in the
synthetic benchmark; it is not introduced into the real-audit comparison after
observing those events.

Predeclared sensitivity analyses vary one factor at a time: donor-screening
rules, 45/60/90-day effect windows, raw versus log reporting scale, and
evidence-tier thresholds. The sensitivity rows reveal availability and
interpretation dependence; they are not a search for the setting with the most
favorable residual pattern. The standard synthetic-control rows in the main
and reliability-ablation experiments are checked for exact alignment to
tolerance $10^{-10}$, preventing a seed or input mismatch from masquerading
as an ablation finding.

\subsection{Interval coverage protocol}

Fixed-weight interval coverage is evaluated on held-out synthetic effects. The
confidence level, seven-day block length, perturbation types, and 1,000
repetitions are fixed in the frozen configuration. The separate full
selection-aware synthetic coverage protocol was frozen before outcomes. Its
feasibility assessment records that full donor re-screening and weight
refitting at the specified nested scale is infeasible within the submission
deadline. The report therefore does not convert real-event selection-aware
intervals into coverage-calibrated physical-bias intervals.

\subsection{External context and reproducibility}

Same-site alternate-POC records, narrow hourly POC windows, QA availability,
and targeted official-document review supply contextual checks with known
limitations. A separate parameter-88502 pipeline uses its own scan and audit;
it is never pooled with 88101. The frozen release records the exact source
commit, locked Python environment, source manifests, machine-readable checks,
and a two-environment comparison of designated core artifact hashes.
